\documentclass[10pt, letterpaper]{article}

% Packages:
\usepackage[
    ignoreheadfoot, % set margins without considering header and footer
    top=2 cm, % seperation between body and page edge from the top
    bottom=2 cm, % seperation between body and page edge from the bottom
    left=2 cm, % seperation between body and page edge from the left
    right=2 cm, % seperation between body and page edge from the right
    footskip=1.0 cm % seperation between body and footer
    % showframe % for debugging 
]{geometry} % for adjusting page geometry
\usepackage{titlesec} % for customizing section titles
\usepackage{tabularx} % for making tables with fixed width columns
\usepackage{array} % tabularx requires this
\usepackage[dvipsnames]{xcolor} % for coloring text
\definecolor{primaryColor}{RGB}{0, 0, 0} % define primary color
\usepackage{enumitem} % for customizing lists
\usepackage{fontawesome5} % for using icons
\usepackage{amsmath} % for math
\usepackage[
    pdftitle={Mustafa Alp Ekici's CV},
    pdfauthor={Mustafa Alp Ekic},
    pdfcreator={LaTeX with RenderCV},
    colorlinks=true,
    urlcolor=primaryColor
]{hyperref} % for links, metadata and bookmarks
\usepackage[pscoord]{eso-pic} % for floating text on the page
\usepackage{calc} % for calculating lengths
\usepackage{bookmark} % for bookmarks
\usepackage{lastpage} % for getting the total number of pages
\usepackage{changepage} % for one column entries (adjustwidth environment)
\usepackage{paracol} % for two and three column entries
\usepackage{ifthen} % for conditional statements
\usepackage{needspace} % for avoiding page brake right after the section title
\usepackage{iftex} % check if engine is pdflatex, xetex or luatex

% Ensure that generate pdf is machine readable/ATS parsable:
\ifPDFTeX
    \input{glyphtounicode}
    \pdfgentounicode=1
    \usepackage[T1]{fontenc}
    \usepackage[utf8]{inputenc}
    \usepackage{lmodern}
\fi

\usepackage{charter}

% Some settings:
\raggedright
\AtBeginEnvironment{adjustwidth}{\partopsep0pt} % remove space before adjustwidth environment
\pagestyle{empty} % no header or footer
\setcounter{secnumdepth}{0} % no section numbering
\setlength{\parindent}{0pt} % no indentation
\setlength{\topskip}{0pt} % no top skip
\setlength{\columnsep}{0.15cm} % set column seperation
\pagenumbering{gobble} % no page numbering

\titleformat{\section}{\needspace{4\baselineskip}\bfseries\large}{}{0pt}{}[\vspace{1pt}\titlerule]

\titlespacing{\section}{-1pt}{0.3 cm}{0.2 cm} % section title spacing

\renewcommand\labelitemi{$\vcenter{\hbox{\small$\bullet$}}$} % custom bullet points

\newenvironment{highlights}{
    \begin{itemize}[
        topsep=0.10 cm,
        parsep=0.10 cm,
        partopsep=0pt,
        itemsep=0pt,
        leftmargin=0 cm + 10pt
    ]
}{
    \end{itemize}
} % new environment for highlights

\newenvironment{highlightsforbulletentries}{
    \begin{itemize}[
        topsep=0.10 cm,
        parsep=0.10 cm,
        partopsep=0pt,
        itemsep=0pt,
        leftmargin=10pt
    ]
}{
    \end{itemize}
} % new environment for highlights for bullet entries

\newenvironment{onecolentry}{
    \begin{adjustwidth}{0 cm + 0.00001 cm}{0 cm + 0.00001 cm}
}{
    \end{adjustwidth}
} % new environment for one column entries

\newenvironment{twocolentry}[2][]{
    \onecolentry
    \def\secondColumn{#2}
    \setcolumnwidth{\fill, 4.5 cm}
    \begin{paracol}{2}
}{
    \switchcolumn \raggedleft \secondColumn
    \end{paracol}
    \endonecolentry
} % new environment for two column entries

\newenvironment{threecolentry}[3][]{
    \onecolentry
    \def\thirdColumn{#3}
    \setcolumnwidth{, \fill, 4.5 cm}
    \begin{paracol}{3}
        {\raggedright #2} \switchcolumn
}{
    \switchcolumn \raggedleft \thirdColumn
    \end{paracol}
    \endonecolentry
} % new environment for three column entries

\newenvironment{header}{
    \setlength{\topsep}{0pt}\par\kern\topsep\centering\linespread{1.5}
}{
    \par\kern\topsep
} % new environment for the header

\newcommand{\placelastupdatedtext}{%
  \AddToShipoutPictureFG*{%
    \put(
        \LenToUnit{\paperwidth-2 cm-0 cm+0.05cm},
        \LenToUnit{\paperheight-1.0 cm}
    ){\vtop{{\null}\makebox[0pt][c]{
        \small\color{gray}\textit{Last updated in September 2024}\hspace{\widthof{Last updated in September 2024}}
    }}}%
  }%
}%

% save the original href command in a new command:
\let\hrefWithoutArrow\href

\begin{document}

\newcommand{\AND}{\unskip
    \cleaders\copy\ANDbox\hskip\wd\ANDbox
    \ignorespaces
}
\newsavebox\ANDbox
\sbox\ANDbox{$|$}

\begin{header}
    \fontsize{25 pt}{25 pt}\selectfont Mustafa Alp Ekici

    \vspace{5 pt}

    \fontsize{9 pt}{10 pt}\selectfont

    \mbox{Milano, Italy}%
    \kern 3.0 pt%
    \AND%
    \kern 3.0 pt%
    \mbox{\hrefWithoutArrow{mailto:mustafalpekici@gmail.com}{mustafalpekici@gmail.com}}%
    \kern 3.0 pt%
    \AND%
    \kern 3.0 pt%
    \mbox{\hrefWithoutArrow{tel:+39-333-829-7495}{+39 333 829 7495}}%
    \kern 3.0 pt%
    \AND%
    \kern 3.0 pt%
    \mbox{\hrefWithoutArrow{https://linkedin.com/in/mustafalpekici/}{linkedin.com/in/mustafalpekici}}%
    \kern 3.0 pt%
    \AND%
    \kern 3.0 pt%
    \mbox{\hrefWithoutArrow{https://mustafalpekici.vercel.app/}{mustafalpekici.vercel.app}}%

\end{header}

\vspace{5 pt - 0.3 cm}

% ======================
% Education
% ======================
\section{Education}


\begin{twocolentry}{Sept 2026 – Present}
    \textbf{Politecnico di Milano}, MSc in Electronics Engineering
\end{twocolentry}

\vspace{0.2 cm}


\begin{twocolentry}{Sept 2022 – June 2026}
    \textbf{Middle East Technical University}, BS in Electrical and Electronics Engineering
\end{twocolentry}

\vspace{0.10 cm}
\begin{onecolentry}
    \begin{highlights}
        \item GPA: 3.61/4.0 (\href{https://drive.google.com/file/d/1neF2DpCgctMv6hOJC8BwWndFErBjKNRt/view?usp=drive_link}{Transcript})
        \item \textbf{Specialization:} Electronics, Biomedical
    \end{highlights}
\end{onecolentry}

\begin{twocolentry}{Sept 2020 – June 2022}
    \textbf{Middle East Technical University}, BS in Computer Engineering
\end{twocolentry}

\vspace{0.10 cm}
\begin{onecolentry}
    \begin{highlights}
        \item Transferred to Electrical and Electronics Engineering Department
    \end{highlights}
\end{onecolentry}

% ======================
% Experience
% ======================
\section{Experience}

\begin{twocolentry}{July 2025 – Sept 2025}
    \textbf{Research Intern ERASMUS+ Traineeship}, TU Delft -- Delft, Netherlands
\end{twocolentry}

\vspace{0.10 cm}
\begin{onecolentry}
    \begin{highlights}
        \item Designed and simulated Love-mode SAW biosensors in COMSOL, analyzed the effect of guiding layer thickness on acoustic wave propagation through eigenfrequency studies.
        \item Developed MATLAB algorithms to simulate biofilm formation and established a LiveLink connection with COMSOL for automated 3D biofilm generation and integration into sensor models.
        \item Optimized sensor geometry and materials in COMSOL and analyzed phase shifts after integrating bacteria, correlating phase variations with bacterial quantity to evaluate sensor performance.
    \end{highlights}
\end{onecolentry}

\vspace{0.2 cm}
\begin{twocolentry}{Sept 2024 – April 2025}
    \textbf{Part-time Engineer}, METU MEMS Center -- Ankara, Turkey
\end{twocolentry}

\vspace{0.10 cm}
\begin{onecolentry}
    \begin{highlights}
        \item Conducted research on delta-sigma ADC architectures, readout circuitry for micro-g MEMS accelerometers, and sensor production methods within the scope of the project.
    \end{highlights}
\end{onecolentry}

\vspace{0.2 cm}
\begin{twocolentry}{August 2024 – Sept 2024}
    \textbf{Engineering Intern}, Avionics Systems Engineering Department, Roketsan -- Ankara, Turkey
\end{twocolentry}

\vspace{0.10 cm}
\begin{onecolentry}

    \begin{highlights}
        \item Gained in-depth knowledge of avionics systems and systems engineering principles, with a focus on their application in avionics.
        \item Conducted power processing analysis using LTspice, designing and simulating DC-DC converters and Pi filters while developing practical skills in electronic circuit simulation and analysis.
    \end{highlights}
\end{onecolentry}

\vspace{0.2 cm}
\begin{twocolentry}{July 2024}
    \textbf{Research Intern}, UMRAM -- Ankara, Turkey
\end{twocolentry}

\vspace{0.10 cm}
\begin{onecolentry}
    \begin{highlights}
        \item Designed and implemented RF Bias Tee circuits and Low Noise Amplifiers for 3T MRI systems using Altium Designer and Proteus, gaining hands-on experience in PCB design, soldering, and testing with network analyzers, while collaborating with Prof. Dr. Ergin Atalar and a team of engineers to optimize RF circuit designs in a leading MRI research center.
    \end{highlights}
    \vspace{0.1cm}
    Full internship report available \href{https://drive.google.com/file/d/1wPOYxnAS0WKuvXy8KdiaDvrH2BX0o-4U/view}{[here]}.
\end{onecolentry}

% ======================
% Projects
% ======================
\section{Projects}

\vspace{0.2 cm}
\begin{twocolentry}{}
    \textbf{Analog IC Design: Operational Amplifier, Bandgap Reference \& LDO}
\end{twocolentry}

\vspace{0.10 cm}
\begin{onecolentry}
    \begin{highlights}
        \item Designed and simulated a two stage CMOS operational amplifier, bandgap reference, and LDO regulator in XFAB 180 nm CMOS using Cadence.
        \item Verified the design across PVT corners, achieving over 97 dB op-amp DC gain, over 5 MHz unity-gain bandwidth, and over 93 dB CMRR while maintaining sub-100 $\mu$A current consumption.
        \item Integrated the complete regulator system, achieving approximately 1.794 V output, 1.4\% bandgap variation, over $60^\circ$ phase margin, and up to $-67$ dB PSRR at 10 kHz.
    \end{highlights}
    Full report available at: \href{https://drive.google.com/file/d/1EyeFLb7r0vsK6oKQAR1TiKlAlgsI0LLn/view?usp=drive_link}{[here]}.
\end{onecolentry}

\vspace{0.2 cm}
\begin{twocolentry}{}
    \textbf{Review on Silicon-Germanium (SiGe) Technology}
\end{twocolentry}

\vspace{0.10 cm}
\begin{onecolentry}
    \begin{highlights}
        \item Wrote a detailed technical review discussing the bandgap engineering and strain physics pertaining to SiGe Heterojunction Bipolar Transistors (HBTs).
        \item Analyzed cutting-edge fabrication methods such as UHV/CVD and deep-trench isolation for devices operating at sub-THz frequencies.
        \item Investigated the reliability of SiGe devices in extreme environments and assessed future trends in sub-THz applications and monolithic optoelectronics.
    \end{highlights}
        Full report available at: \href{https://drive.google.com/file/d/11weppkvTCNk7jUJ_5FexN0yiupOhWysw/view?usp=sharing}{[here]}.
\end{onecolentry}

\vspace{0.2 cm}
\begin{twocolentry}{}
    \textbf{Neural-Network MAC Tile Accelerator Design (VLSI)}
\end{twocolentry}

\vspace{0.10 cm}


\begin{onecolentry}
    \begin{highlights}
        \item Designed a Neural Network Accelerator on XFAB 180nm completing the RTL to Physical Layout flow with Cadence Genus  \&  Innovus.
        \item Developed a Dual-MAC architecture achieving 2x throughput and 44\% energy reduction compared to the Single-MAC design.
        \item Achieved timing closure at 50 MHz and completed physical verification (DRC/LVS).
        
    \end{highlights}
\end{onecolentry}
    Full report available at: \href{https://drive.google.com/file/d/1bxllJ0fdsJWdAEdpuMTNm5sW8nCLTTSL/view?usp=sharing}{[here]}.

\vspace{0.2 cm}
\begin{twocolentry}{}
    \textbf{X-Ray CT Simulation \& Image Reconstruction}
\end{twocolentry}

\vspace{0.10 cm}
\begin{onecolentry}
    \begin{highlights}
        \item Developed from scratch in MATLAB the Forward Projection and Inverse Reconstruction algorithms.
        \item Developed the Ray-Driven Exact Path Length method for accurate sinogram creation and performed sensitivity analysis on sampling parameters.
        \item Demonstrated that Filtered Backprojection decreased reconstruction error by around 80\% when compared to Unfiltered Backprojection.
    \end{highlights}
        Full report available at: \href{https://drive.google.com/file/d/1aaxeCN1S0WdStdJdRa9rEZZdUCy7QbqD/view?usp=sharing}{[here]}.
\end{onecolentry}


\vspace{0.2 cm}

\begin{twocolentry}{}
    \textbf{Micro Air Conditioner Designing}
\end{twocolentry}

\vspace{0.10 cm}
\begin{onecolentry}
    \begin{highlights}
        \item Designed and implemented an analog micro air conditioner by integrating sensing, control unit, operation unit and RGB LED display; performed LTspice simulations and breadboard prototyping to realize the design.
        \item Conducted simulations and experimental tests to achieve autonomous temperature regulation with less than $\pm 0.8^\circ$C error, including validation of heating and cooling performance 
        \item  Programs Used: LTspice
    \end{highlights}
    Full report available at: \href{https://drive.google.com/file/d/14LxSZs8bfI0C_sf_S6rY0oO5PZYFHFg2/view}{[here]}.
\end{onecolentry}

\vspace{0.2 cm}
\begin{twocolentry}{}
    \textbf{Power Cable Selection Interface Design}
\end{twocolentry}

\vspace{0.10 cm}
\begin{onecolentry}
    \begin{highlights}
        \item Designed and developed a Python-based GUI tool to automate power cable selection by integrating electrical engineering calculations with an intuitive interface.
        \item Implemented algorithms to calculate current rating, voltage drop, line losses, voltage regulation, and applied temperature and trench correction factors based on international standards.
        \item Performed a 10-year economic analysis combining installation costs with long-term energy loss expenses to optimize cost efficiency.
        \item Programs Used: Python (PyQt5, Pandas)
    \end{highlights}
    Full report available at: \href{https://drive.google.com/file/d/16f8ijAeEyLVm5ekeazxyMJRlgb9_2rDH/view}{[here]}.
\end{onecolentry}



% ======================
% Technologies
% ======================
% ======================
% Skills
% ======================
\section{Skills}

\begin{onecolentry}
    \textbf{Programming \& Software:} Sentaurus TCAD, Cadence, Python (PyQt5, Pandas), MATLAB, COMSOL Multiphysics (LiveLink), LTspice, Altium Designer, Proteus, KiCad
\end{onecolentry}


\section{Languages}

\begin{onecolentry}
  \textbf{Turkish:} Native
\end{onecolentry}

\vspace{0.15 cm}
\begin{onecolentry}
    \textbf{English:} C2
\end{onecolentry}

\vspace{0.15 cm}
\begin{onecolentry}
    \textbf{German:} A2
\end{onecolentry}

% ======================
% Certifications
% ======================
\section{Certifications}

\begin{onecolentry}
    \textbf{Cleanroom Training Certificate — METU MEMS Center}\\
    Trained and certified for cleanroom work, gaining practical experience with entry procedures, safety rules, and proper dressing 
\end{onecolentry}

% ======================
% Honors & Awards
% ======================



\end{document}
